
\[
\begin{aligned}
\lambda'_u(1+Xt) &=\Ghat(\rho_t)(\lambda_u(\rho_t^{-1}(1+Xt)))\\
&=\Ghat(\rho_t)\left(\lambda_u\left(\frac{1+(X+1)t}{1+t}\right)\right)\\
&=\Ghat({}_t)\left(\frac{\underline{1}_o+\{1+(X+1)t\}u}{\underline{1}_o+\{1+t\}u}\right)\\
&=\underline{1}_\nabla+\frac{\{1+Xt\}u}{\underline{1}_\nabla+\underline{1}_\nabla.u}\\
&=\rho_u^{-1}(\underline{1}_\nabla+\{1+Xt\}u).
\end{aligned}
\]

Si on note \(\gamma'\) les opérations \(\gamma\) correspondant aux opérations \(\lambda'\), on a donc :
\begin{equation}\tag{3.6.2}
  \gamma'_u(1+Xt)=\underline{1}_\nabla+\{1+Xt\}u .
\end{equation}

On notera \(Q'_{i,j}\) le coefficient de \(t^j\) dans \(\gamma'^i(1+\ldots+X_nt^n+\ldots)\)~; le lecteur pourra vérifier que c'est un polynôme de \(Z[X_1,\ldots,X_{ij}]\), dont les composantes isobares sont de poids variant entre \(j\) et \(ij\).

Soient \(K\) un \(\lambda\)-anneau, \(x\in K\). Du diagramme (3.6.1), on déduit les relations :
\begin{equation}\tag{3.6.3}
  \forall i\in\mathbb N,\quad\gamma_t(\lambda^i(x))=\lambda'^i(\gamma_t(x))~;
  \qquad
  \gamma_t(\gamma^i(x))=\gamma'^i(\gamma_t(x)).
\end{equation}
\begin{equation}\tag{3.6.4}
  \gamma^j(\gamma^i(x))
  =Q'_{i,j}(\gamma^1(x),\ldots,\gamma^{ij}(x)).
\end{equation}

On remarquera enfin que la structure de pré-\(\lambda\)-anneau qui vient d'être définie sur \(\Ghat\) est une structure de \(\lambda\)-anneau, puisque déduite de celle qui a été définie en 2.3 par l'isomorphisme \(\rho\).

3.7. Pour éviter les confusions, le foncteur \(\Ghat\), muni de la structure de foncteur en \(\lambda\)-anneaux définie au §2 (resp. au §3) sera noté \(\Ghat_o\) (resp. \(\Gnabla\)). Il résulte des considérations précédentes que pour se donner une
